\section{实践案例与落地路径}

\subsection{案例：Multi-hop DSPy 系统}

一个典型的 multi-hop pipeline 先检索、再推理、然后验证，再调用工具。DSPy 能让你把这几个步骤分别抽象成 module，再用 optimizer 帮你搜索 prompt/示例组合。实践中会遇到 `chain-of-thought` 与 `tool invocation` 在时间轴上的同步问题。

\subsection{治理与文档}

推荐在 repo 中新增 `DSPy.md` 文档，包含 deployment 版本、Action Timeline、Evidence Matrix 链接以及 security checklist。每次 release 将 `DSPy.md` 作为 release note 的一部分并交由 QA 检查。

\begin{importantbox}{文档就是治理}
DSPy 代码比普通 pipeline 多出两层：DSL + optimizer log。如果不写文档，开发者之间就会出现 ``我不知道 optimizer 为什么改了 prompt'' 的摩擦。文档应包含：signature、metric、fallback 条件、owner 联系方式。
\end{importantbox}

\subsection{本章小结}

把大型复合系统拆解成可复用案例，并为其维持明文文档，是从概念验证走向可运营的关键一步。

